% Literature review. \citeneeded{...} marks claims whose reference the team still has to add
% (they render in red, so none can slip into the submitted version).
\chapter{Literature review}
\label{ch:literature}

This chapter reviews the four bodies of work the thesis connects: tail-risk measures and their estimation
(\cref{sec:lit-risk}), jump-diffusion models and their numerical solution (\cref{sec:lit-jumps}),
physics-informed and other neural PDE solvers (\cref{sec:lit-pinn}), and their use in finance
(\cref{sec:lit-pinnfin}). \Cref{sec:lit-gap} states the gap.

% ------------------------------------------------------------------------------------------------
\section{Tail-risk measures and their estimation}
\label{sec:lit-risk}

\paragraph{VaR and ES.} Value at Risk became the industry standard in the 1990s \citep{jorion2007}. It is a
quantile, and therefore says nothing about the size of losses beyond it; it is also not subadditive, so the VaR of
a portfolio can exceed the sum of its parts. \Citet{artzner1999} defined the axioms of coherent risk measures, and
\citet{acerbi2002} showed that Expected Shortfall satisfies them. The Basel Committee's Fundamental Review of the
Trading Book replaced 99\% VaR with 97.5\% ES for market-risk capital \citep{bcbs2019}.

\paragraph{Estimation methods.} Four families dominate practice.
\begin{itemize}
  \item \textbf{Parametric} models assume a distribution, usually Gaussian or Student-$t$, with volatility from an
    EWMA filter \citep{riskmetrics1996} or a GARCH model \citep{bollerslev1986}. Multivariate versions model
    correlations, for example with dynamic conditional correlation \citep{engle2002}.
  \item \textbf{Historical simulation} uses the empirical distribution of past returns. It makes no distributional
    assumption but reacts slowly to changes in volatility. Filtered historical simulation rescales past returns by
    the ratio of current to past volatility \citep{barone1999}, and is a strong benchmark in comparative studies
    \citeneeded{Bank of England study of filtered-HS VaR, cited in Review 2}.
  \item \textbf{Extreme value theory} fits a generalised Pareto distribution to the tail. \Citet{mcneil2000}
    combine it with GARCH filtering, applying EVT to the standardised residuals.
  \item \textbf{Monte Carlo} simulation from a fitted model handles any model and any portfolio, including
    options, at a cost that grows with the precision required in the tail \citep{glasserman2003}.
\end{itemize}

\paragraph{Backtesting.} \Citet{kupiec1995} tests whether the frequency of VaR breaches matches the nominal
level, and \citet{christoffersen1998} adds a test that breaches are independent over time. Neither evaluates ES.
ES is not elicitable on its own \citep{gneiting2011}, which long made it hard to backtest, but the pair (VaR, ES)
is jointly elicitable \citep{fissler2016}. This gives strictly consistent scoring functions such as the FZ0 loss
of \citet{patton2019}, with which competing forecasts can be ranked, and their differences tested with the
Diebold--Mariano test \citep{diebold1995} using a Newey--West variance \citep{newey1987}.

% ------------------------------------------------------------------------------------------------
\section{Jump-diffusion models}
\label{sec:lit-jumps}

Returns exhibit fat tails, volatility clustering and occasional discontinuities \citep{mandelbrot1963,cont2001}.
\Citet{merton1976} added a compound Poisson process with Gaussian jump sizes to the Black--Scholes diffusion; the
density is then a Poisson mixture of Gaussians, available in closed form. \Citet{kou2002} uses double-exponential
jumps, which give a heavier tail and analytical tractability for path-dependent options. \Citet{cont2004} give a
comprehensive treatment of jump processes in finance, including their generators and the integro-differential
equations they lead to.

The density of a jump-diffusion solves a forward PIDE, the Fokker--Planck equation with an added nonlocal jump
term. It is usually solved by finite differences, treating the local terms implicitly and the jump integral
explicitly, with the convolution computed by FFT \citep{andersen2000,cont2004}. These schemes are accurate and fast
in one dimension, but as with every grid method their cost grows as $n^d$ in $d$ dimensions, which limits them to
two or three assets.

% ------------------------------------------------------------------------------------------------
\section{Neural solvers for PDEs}
\label{sec:lit-pinn}

\paragraph{Physics-informed neural networks.} The idea of training a network to satisfy a differential equation
goes back to \citet{lagaris1998}. \Citet{raissi2019} gave it its current form: a network approximates the solution,
automatic differentiation evaluates the PDE residual at sampled collocation points, and the loss combines the
residual with initial and boundary conditions. PINNs are mesh-free, handle forward and inverse problems in the
same framework and can take parameters as inputs \citep{karniadakis2021}.

\paragraph{Known difficulties.} PINN training is fragile. Loss terms can have very different gradient magnitudes,
which unbalances training \citep{wang2021}, and some problems are not learned at all because of the optimisation
landscape rather than the network's capacity \citep{krishnapriyan2021}. Two difficulties are specific to densities:
a density spans many orders of magnitude, so a residual on $p$ ignores the tail, and the PDE is unchanged by
scaling $p$, so normalisation must be imposed separately.

\paragraph{Fokker--Planck equations.} PINNs have been applied to forward and inverse Fokker--Planck problems
\citep{chen2021fp}, mostly with continuous (diffusive) dynamics. Nonlocal jump terms, which make the equation a
PIDE, are less studied.

\paragraph{High dimensions.} Deep BSDE methods \citep{han2018} and the deep Galerkin method
\citep{sirignano2018} solve PDEs in hundreds of dimensions, typically for a value at one point or for a smooth
solution, not for a density across its tail. In PINNs the cost of second derivatives grows with $d$; Hutchinson's
stochastic trace estimator \citep{hutchinson1989} replaces the full Laplacian by random projections, and has been
used to scale PINNs to high dimensions \citep{hu2024}.

\paragraph{Neural operators and density models.} Neural operators such as DeepONet \citep{lu2021} and the Fourier
neural operator \citep{li2021} learn a map from parameters or initial conditions to solutions, from many solved
examples. The amortised PINN of this thesis targets the same goal, one network for a family of problems, but is
trained from the equation alone, without solved examples. Data-driven density estimators such as normalising
flows \citep{papamakarios2021} learn densities from samples alone; the data-only baseline of this thesis is the
simplest such model, built from the same network as the PINN so that the comparison isolates the physics.

% ------------------------------------------------------------------------------------------------
\section{Neural networks for finance}
\label{sec:lit-pinnfin}

Option pricing has been the main financial application of neural PDE solvers: the deep Galerkin method was
demonstrated on American options \citep{sirignano2018}, deep BSDE methods on high-dimensional pricing
\citep{han2018}, and PINNs on the Black--Scholes equation and on jump-diffusion pricing PIDEs
\citeneeded{PINN option pricing, incl.\ jump-diffusion, cited in Review 2}. These works solve the backward
pricing equation for a price, a smooth function whose accuracy is judged in the bulk. Tail-risk work with neural
networks has mostly been data-driven, using networks to forecast VaR or ES directly from past returns
\citeneeded{neural VaR/ES forecasting, e.g.\ quantile-regression networks}.

% ------------------------------------------------------------------------------------------------
\section{Research gap}
\label{sec:lit-gap}

Applying a PINN to a financial equation is no longer new. What the literature does not provide is a systematic
test of whether PINN-solved \emph{densities} give accurate and useful \emph{tail risk}. Such a test requires:
\begin{itemize}
  \item tail accuracy measured separately from bulk accuracy, against an exact reference;
  \item VaR and ES compared with both numerical solvers (finite differences, Monte Carlo) and the statistical risk
    models used in practice (historical simulation, GARCH, EVT), including joint VaR--ES scoring;
  \item the multi-asset case, where grid solvers cannot be used and a neural solver could add value;
  \item an honest accounting of training cost and query latency against Monte Carlo;
  \item a direct measurement of what the physics contributes over the same network trained on data, and an
    account of where the method fails.
\end{itemize}
This thesis provides that test, for jump-diffusion models with up to five assets and on two real equity baskets.
